\chapter{The Interstellar Medium and Gravitational Collapse}

Chapters~1--7 treated a star (or a pulsar, or a binary) essentially as an isolated point source: a clock, a moving mass, a blackbody, a photon-emitting aperture problem. In reality, no photon reaches a telescope without first traversing tens to thousands of parsecs of gas and dust --- the \emph{interstellar medium} (ISM). The ISM scatters, absorbs, and disperses the light we use for every measurement in this book, and it is also the raw material out of which every star is born and into which every star eventually returns its processed mass. This chapter first characterizes the ISM as an absorbing, scattering, and line-emitting medium, and then turns to its own internal dynamics: the virial theorem, the Jeans criterion for gravitational instability, and the free-fall collapse that marks the transition from a diffuse cloud to a forming protostar.

\section{Composition, Density, and Temperature}

\begin{definition}[Interstellar Medium]
	The \emph{interstellar medium} is the matter --- gas, dust, cosmic rays, and magnetic field --- occupying the space between stars within a galaxy.
\end{definition}

Despite containing roughly $20\%$ of the ordinary (baryonic) mass of the Milky Way, the ISM is extraordinarily dilute:

\begin{table}[htbp]
	\centering
	\begin{tabular}{@{}l l@{}}
		\toprule
		\textbf{Property} & \textbf{Typical Value} \\
		\midrule
		Number density $n$ & $10^{-4} \text{--} 10^{6}\,\mathrm{cm^{-3}}$ \\
		Temperature $T$ & $10\text{--}10^{6}\,\mathrm{K}$ \\
		Composition (by mass) & $70\%$ H, $28\%$ He, $\lesssim 2\%$ heavier elements \\
		Phase split (by mass) & $99\%$ gas, $1\%$ dust \\
		Magnetic field strength & $\sim 10^{-10}\,\mathrm{T}$ ($\sim 1\,\mu\mathrm{G}$) \\
		\bottomrule
	\end{tabular}
	\caption{Characteristic bulk properties of the interstellar medium.}
	\label{tab:ism_bulk_properties}
\end{table}

For comparison, the best vacuum achievable in a terrestrial laboratory (extreme high vacuum, XHV) reaches $\sim 10^{5}\,\mathrm{cm^{-3}}$: the diffuse phase of the ISM is routinely \emph{more rarefied than our best laboratory vacuum}. This has an important physical consequence exploited throughout this chapter: at such densities, collisional de-excitation is so infrequent that highly forbidden atomic and molecular transitions --- undetectable in any terrestrial spectrometer --- become the dominant, and often only, radiative channel available to interstellar gas (Section~\ref{sec:ism_hydrogen}).

The gas phase consists of ions (e.g.\ $\mathrm{H^+}$, $\mathrm{He^+}$, $\mathrm{He^{2+}}$), free electrons, and molecules (e.g.\ $\mathrm{H_2}$, CO, CN, $\mathrm{H_2O}$, $\mathrm{NH_3}$). The dust phase consists of sub-micron grains ($\sim 1\,\mathrm{nm}$ to $\sim 1\,\mu\mathrm{m}$) of silicates, amorphous carbon, and metals (Fe, Ni), often mantled with ices ($\mathrm{H_2O, NH_3}$). Although dust carries only $1\%$ of the ISM's mass, its large geometric cross-section per unit mass makes it the dominant agent of interstellar extinction (Section~\ref{sec:ism_extinction}) and the primary catalyst for $\mathrm{H_2}$ formation (Section~\ref{sec:ism_hydrogen}).

Studying the ISM matters for three independent reasons: (i) every photon used in Chapters 1--7 is filtered by it, so its optical properties must be understood before extinction and reddening corrections (Chapter~4) can be trusted; (ii) it hosts physical and chemical regimes --- extreme rarefaction, extreme cold, non-equilibrium chemistry --- inaccessible to any laboratory; and (iii) it is both the birthplace and the graveyard of stars, which is the subject of the second half of this chapter.

\section{Radiative Interaction with Dust: Extinction, Scattering, and Reddening}
\label{sec:ism_extinction}

\begin{definition}[Extinction]
	\emph{Extinction} is the total attenuation of light passing through a medium, defined as the sum of true absorption and scattering out of the line of sight:
	\begin{equation}
		\text{Extinction} = \text{Absorption} + \text{Scattering}.
		\label{eq:extinction_decomposition}
	\end{equation}
\end{definition}

Absorption occurs either as sharp line absorption by atoms and molecules, or as broadband continuum absorption by larger dust grains. Scattering off a grain of radius $a$ falls into three regimes depending on the ratio of grain size to wavelength:

\begin{table}[htbp]
	\centering
	\begin{tabular}{@{}l l l l@{}}
		\toprule
		\textbf{Regime} & \textbf{Condition} & \textbf{Cross-section scaling} & \textbf{Physical origin} \\
		\midrule
		Geometric & $a \gg \lambda$ & $\sigma \propto a^2$ & Ray optics; grain blocks its own geometric shadow. \\
		Mie & $a \sim \lambda$ & $\sigma \propto a^3\lambda^{-1}$ & Full electromagnetic (Mie) solution; no closed form. \\
		Rayleigh & $a \ll \lambda$ & $\sigma \propto a^6\lambda^{-4}$ & Grain acts as a driven electric dipole. \\
		\bottomrule
	\end{tabular}
	\caption{Scattering regimes for a spherical grain of radius $a$ illuminated at wavelength $\lambda$.}
	\label{tab:scattering_regimes}
\end{table}

\paragraph{Physical origin of the three regimes.} The geometric limit is simple ray optics: a grain much larger than the wavelength blocks light over its own projected area, $\sigma = \pi a^2$, independent of $\lambda$. The Rayleigh limit is a classical driven-dipole problem: an incident wave of amplitude $E_0$ induces a dipole moment $p = \alpha E_0$ in the grain, with polarizability $\alpha \propto a^3$ (the response of a sub-wavelength dielectric or conducting sphere scales with its volume). This oscillating dipole radiates according to the Larmor formula, $P_{\rm rad} \propto \ddot p^{\,2} \propto \omega^4 p^2$, so the re-radiated power scales as $\omega^4 a^6 E_0^2 \propto a^6/\lambda^4$ (using $\omega = 2\pi c/\lambda$); dividing by the incident intensity $I \propto E_0^2$ gives the scattering cross-section
\begin{keyresult}
	\begin{equation}
		\sigma_{\rm Rayleigh} \propto \frac{a^6}{\lambda^4}.
		\label{eq:rayleigh_scattering_cross_section}
	\end{equation}
\end{keyresult}
This is precisely the same dipole-scattering result invoked in Chapter~5 (Section~\ref{sec:windows}) for air molecules rather than dust grains, and is the microscopic origin of the $\lambda^{-4}$ law responsible for the blue daytime sky. The intermediate Mie regime, by contrast, has no elementary derivation of this kind: the exact solution requires expanding the scattered field in an infinite series of vector spherical harmonics (the Mie series), whose terms do not collapse to a simple power law; the scaling $\sigma \propto a^3\lambda^{-1}$ quoted in Table~\ref{tab:scattering_regimes} is only the effective average behaviour across this transition regime, not a closed-form result.

Interstellar grains ($a \sim 0.01\text{--}1\,\mu\mathrm{m}$) are comparable to or smaller than optical wavelengths, so ISM scattering is generically in the Mie or Rayleigh regime, both of which share the qualitative feature that $\sigma$ \emph{decreases} with increasing $\lambda$. A beam of white starlight passing through a dusty cloud therefore loses blue photons preferentially, emerging both dimmer (extinction) and redder (\emph{interstellar reddening}) than it left the source. This is the microphysical origin of the empirical extinction law $A_V = R_V\,E(B-V)$ with $R_V \approx 3.1$ already introduced in Chapter~4 (Eq.~\eqref{eq:rv_extinction}): the observed $\gamma \sim 1$ power-law is the net effect of averaging over a realistic grain-size distribution straddling the Mie and Rayleigh regimes, rather than a single clean $\lambda^{-4}$ Rayleigh law. A well-known consequence is visible directly in absorption and reflection nebulae: dense dust lanes such as Barnard 68 appear as near-total silhouettes in the optical yet become largely transparent in the infrared (where $a \ll \lambda$ deep in the Rayleigh regime and $\sigma$ is small), while starlight scattered off the dust surrounding the Pleiades produces a blue reflection nebula for the same Rayleigh-type reason that the terrestrial daytime sky is blue.

\section{Spectral Lines: Formation and Broadening}
\label{sec:ism_lines}

Atoms and ions possess only electronic transitions; molecules additionally possess rotational and vibrational transitions, giving them a far richer spectrum. An excited species relaxing spontaneously produces an \emph{emission} line; a cooler foreground species illuminated by a hotter background continuum source produces an \emph{absorption} line. In neither case is the line infinitely sharp — three independent physical mechanisms set its width.

\subsection{Natural Broadening}

By the time-energy uncertainty relation, a state with a finite lifetime $\Delta\tau$ cannot have an arbitrarily sharp energy:
\begin{equation}
	\Delta E\,\Delta\tau \sim \hbar.
	\label{eq:natural_broadening_energy}
\end{equation}
Since $E = h\nu$, this maps directly onto a minimum frequency width
\begin{equation}
	\Delta\nu_{\rm nat} \sim \frac{1}{2\pi\,\Delta\tau},
	\label{eq:natural_broadening_freq}
\end{equation}
where $\Delta\tau$ is the mean radiative lifetime of the upper state against spontaneous decay. Short-lived (strongly allowed) transitions are intrinsically broad; long-lived (forbidden) transitions are intrinsically narrow — a point revisited for the $21\,\mathrm{cm}$ line below.

\paragraph{Shape: the Lorentzian profile.} The uncertainty argument above fixes only the \emph{width}; the line's actual \emph{shape} follows from modelling the radiating atom classically as a damped oscillator whose emitted wave train decays exponentially once excited. If the radiated energy (intensity $\propto |E(t)|^2$) decays as $e^{-t/\Delta\tau}$, the field amplitude itself decays at half that rate:
\begin{equation}
	E(t) = E_0\,e^{-t/(2\Delta\tau)}\,e^{-i\omega_0 t}\,\Theta(t),
\end{equation}
where $\Theta(t)$ is the step function (the wave train starts at $t=0$, when the atom is excited). The observed spectrum is the power spectrum of this finite wave train, obtained via its Fourier transform:
\begin{equation}
	\tilde E(\omega) = \int_0^\infty E_0\,e^{-t/(2\Delta\tau)}\,e^{-i\omega_0t}\,e^{i\omega t}\,dt = \frac{E_0}{\dfrac{1}{2\Delta\tau} - i(\omega-\omega_0)}.
\end{equation}
Squaring to get the intensity spectrum, $I(\omega) \propto |\tilde E(\omega)|^2$, gives the \textbf{Lorentzian profile}:
\begin{keyresult}
	\begin{equation}
		I(\omega) \propto \frac{1}{(\omega-\omega_0)^2 + \left(\dfrac{1}{2\Delta\tau}\right)^2},
		\label{eq:natural_lorentzian_profile}
	\end{equation}
\end{keyresult}
whose full width at half maximum, $\Delta\omega = 1/\Delta\tau$ (i.e.\ $\Delta\nu=1/(2\pi\Delta\tau)$), exactly reproduces the uncertainty-principle estimate of Eq.~\eqref{eq:natural_broadening_freq}. The defining feature of a Lorentzian --- slowly-decaying power-law wings, $I\propto(\omega-\omega_0)^{-2}$ far from line centre --- is the spectral signature of this exponential (i.e.\ single-exponential-decay) damping process.

\subsection{Pressure (Collisional) Broadening}

Collisions with neighbouring particles interrupt the phase of an emitting atom before it would otherwise have finished radiating, effectively truncating its coherence time. Formally this is identical to natural broadening with $\Delta\tau$ replaced by the mean time between collisions, $\tau_{\rm coll} \sim 1/(n\sigma_{\rm coll} v)$: denser, hotter gas produces broader, more pressure-sensitive lines. Because the underlying physical picture --- an exponentially-truncated wave train --- is identical to the natural-broadening case above (only the truncation mechanism differs: spontaneous decay versus a dephasing collision), the same Fourier-transform argument applies unchanged, and pressure broadening also produces a \textbf{Lorentzian} line shape, Eq.~\eqref{eq:natural_lorentzian_profile}, but with $\Delta\tau \to \tau_{\rm coll}$.

\subsection{Doppler Broadening}

A line-emitting particle moving with line-of-sight velocity $v_\parallel$ has its rest-frame frequency $\nu_0$ Doppler-shifted according to the classical relation already used in Chapter~3 (Eq.~\eqref{eq:doppler_classical}), $\Delta\nu/\nu_0 = v_\parallel/c$. In a gas at temperature $T$, the equipartition theorem assigns mean kinetic energy $\frac{1}{2}k_BT$ to each translational degree of freedom, so the line-of-sight velocity component satisfies $\frac{1}{2}m\avg{v_\parallel^2} = \frac{1}{2}k_BT$, i.e.\ the line-of-sight velocities of individual emitters follow a one-dimensional Maxwell--Boltzmann distribution,
\begin{equation}
	f(v_\parallel)\,dv_\parallel \propto \exp\!\left(-\frac{v_\parallel^2}{2\sigma_v^2}\right)dv_\parallel, \qquad \sigma_v \equiv \sqrt{\avg{v_\parallel^2}} = \sqrt{\frac{k_BT}{m}} \quad (m\text{ the particle mass}).
	\label{eq:maxwell_boltzmann_1d}
\end{equation}

\paragraph{Shape: the Gaussian profile.} Unlike natural or pressure broadening, Doppler broadening is not an intrinsic property of any single emitter --- it is the observational blending of many individually narrow lines, each merely Doppler-shifted by a different emitter's own $v_\parallel$. The line shape is therefore just the velocity distribution itself, rewritten in frequency units. Inverting the classical Doppler relation, $v_\parallel = c(\nu-\nu_0)/\nu_0$, and substituting into Eq.~\eqref{eq:maxwell_boltzmann_1d}:
\begin{equation}
	I(\nu)\,d\nu \propto \exp\!\left[-\frac{c^2(\nu-\nu_0)^2}{2\nu_0^2\sigma_v^2}\right]d\nu,
\end{equation}
which is precisely a \textbf{Gaussian profile} in frequency,
\begin{keyresult}
	\begin{equation}
		I(\nu) \propto \exp\!\left[-\frac{(\nu-\nu_0)^2}{2\sigma_D^2}\right], \qquad \sigma_D = \frac{\nu_0\sigma_v}{c},
		\label{eq:doppler_gaussian_profile}
	\end{equation}
\end{keyresult}
of characteristic (thermal) width
\begin{equation}
	\frac{\Delta\nu_{\rm D}}{\nu_0} \sim \frac{\sigma_v}{c} = \sqrt{\frac{k_B T}{mc^2}}.
	\label{eq:thermal_doppler_width}
\end{equation}
The Gaussian shape is thus a direct fingerprint of the underlying Maxwellian velocity distribution --- in sharp contrast to the power-law-winged Lorentzian of natural/pressure broadening, a Gaussian falls off exponentially fast in its wings, since the population of atoms with any given $|v_\parallel|$ itself falls off exponentially in $v_\parallel^2$. Bulk, non-thermal gas motions (turbulence, infall, rotation) add an additional, generally larger, Doppler contribution with the same Gaussian functional form but $\sigma_v$ replaced by the turbulent or bulk velocity dispersion; a net systemic line-of-sight velocity (e.g.\ Galactic rotation) simply shifts the line centroid $\nu_0$ rather than broadening it.

\paragraph{Remark: the observed profile is a Voigt profile.} A real interstellar line is simultaneously subject to (Lorentzian) natural/pressure broadening and (Gaussian) Doppler broadening; the observed profile is the convolution of the two, known as a \emph{Voigt profile}. In the ISM, Doppler broadening from the low densities and non-thermal turbulence discussed in Section~\ref{sec:ism_heating_cooling} generally dominates over the intrinsically much narrower natural/pressure Lorentzian core, except far out in the line wings, where the slowly-decaying Lorentzian tail can still exceed the exponentially-suppressed Gaussian.

\begin{figure}[htbp]
	\centering
	\begin{tikzpicture}[
		scale=1.15,
		declare function={
			lorentzian(\x) = 1/(1+(\x)^2);
			gaussian(\x) = exp(-ln(2)*(\x)^2);
		}
	]
		\draw[->, thick] (-4.7,0) -- (4.9,0) node[below right, font=\small] {$(\nu-\nu_0)\,/\,(\Delta\nu/2)$};
		\draw[->, thick] (0,-0.15) -- (0,1.65) node[above, font=\small] {$I(\nu)/I_0$};

		\draw[dashed, gray!50, thin] (-4.5,0.5) -- (4.5,0.5);
		\node[gray!70, font=\scriptsize, anchor=east] at (-0.15,0.5) {$0.5$};
		\node[gray!70, font=\scriptsize, anchor=north] at (0,-0.02) {$0$};

		\draw[very thick, accent, domain=-4.5:4.5, samples=200] plot (\x, {lorentzian(\x)});
		\draw[very thick, astroblue, dashed, domain=-4.5:4.5, samples=200] plot (\x, {gaussian(\x)});

		\node[astroblue, font=\small\bfseries, align=center] at (-2.55,1.15) {Gaussian\\\footnotesize(Doppler)};
		\node[accent, font=\small\bfseries, align=center] at (3.15,0.95) {Lorentzian\\\footnotesize(natural / pressure)};
	\end{tikzpicture}
	\caption{The two elementary line-broadening shapes, normalized to the same peak intensity $I_0$ and the same half-width $\Delta\nu/2$. The Lorentzian profile (Eq.~\eqref{eq:natural_lorentzian_profile}) retains extended power-law wings, $I\propto(\nu-\nu_0)^{-2}$, far from line centre, while the Gaussian profile (Eq.~\eqref{eq:doppler_gaussian_profile}) is already negligible by $3$--$4$ half-widths out. The observed \emph{Voigt} profile --- their convolution --- looks essentially Gaussian in its core (where Doppler broadening dominates) but reverts to the Lorentzian's power-law decay far out in the wings.}
	\label{fig:lorentzian_gaussian_comparison}
\end{figure}

\section{Hydrogen in the ISM: HI, HII, \texorpdfstring{H$_2$}{H2}, and the 21 cm Line}
\label{sec:ism_hydrogen}

Hydrogen, the most abundant element in the ISM, exists in three forms: neutral atomic hydrogen ($\mathrm{HI}$), ionized hydrogen ($\mathrm{HII} \equiv \mathrm{H^+}$), and molecular hydrogen ($\mathrm{H_2}$). $\mathrm{HII}$ forms by photoionization of $\mathrm{HI}$ by ultraviolet photons from hot, massive stars. $\mathrm{H_2}$ forms almost exclusively on the surfaces of dust grains: an $\mathrm{HI}$ atom adsorbs onto a grain (raising the probability of a second atom encountering it), and the recombination $2\mathrm{H} \to \mathrm{H_2}$ is sufficiently exothermic that the grain itself must absorb the excess binding energy for the molecule to remain bound (gas-phase formation is far too slow at ISM densities to compete).

\paragraph{The 21 cm Line.} Almost all $\mathrm{HI}$ in the Universe sits in its electronic ground state, and the allowed (Lyman-series) transitions out of that state lie in the ultraviolet, requiring hot background sources to excite them in absorption. Neutral hydrogen instead reveals itself through a completely different, much weaker channel: the \emph{hyperfine transition}, a flip of the relative spin orientation of the proton and electron. This transition is so strongly forbidden that its natural radiative lifetime is $\Delta\tau \sim 10^{7}\,\mathrm{yr}$ — by Eq.~\eqref{eq:natural_broadening_freq}, an extraordinarily narrow natural linewidth, but more importantly, so slow a spontaneous decay rate that under any laboratory density collisional de-excitation (Section~\ref{sec:ism_lines}) always wins first, and the line is never seen in a terrestrial spectrometer. In the diffuse ISM, however, densities are low enough ($n \lesssim 1\,\mathrm{cm^{-3}}$ for HI clouds) that collisions are in fact rare compared to $10^7$ years, and the sheer number of hydrogen atoms along a galactic sightline makes the $21\,\mathrm{cm}$ line easily detectable in emission. Because the line's rest frequency ($1420.4\,\mathrm{MHz}$) is known to extreme precision, its observed Doppler shift (Eq.~\eqref{eq:doppler_classical}) directly maps the line-of-sight velocity of HI gas at every point along the sightline — the single most important observational tool for mapping the distribution and rotation of neutral gas across the Milky Way and other galaxies.

\paragraph{Molecular Hydrogen.} $\mathrm{H_2}$ survives only in regions shielded from dissociating UV radiation by dust (Section~\ref{sec:ism_extinction}); an unshielded molecule is rapidly photodissociated. Being a homonuclear diatomic molecule, $\mathrm{H_2}$ has no permanent dipole moment and therefore no allowed dipole rotational or vibrational transitions, making it extremely difficult to observe directly. Astronomers instead infer its presence and column density using \emph{tracer molecules} such as $\mathrm{CO}$ (which does have a dipole moment), assuming a fixed $\mathrm{H_2}$-to-tracer abundance ratio calibrated independently.

\section{Interstellar Clouds and Nebulae: A Phase Taxonomy}

Physical conditions in the ISM span many orders of magnitude in density and temperature, with no sharp boundaries between regimes; Table~\ref{tab:ism_cloud_types} summarizes the loosely-defined classes.

\begin{table}[htbp]
	\centering
	\small
	\begin{tabular}{@{}l l l l l@{}}
		\toprule
		\textbf{Cloud Type} & \textbf{$T$ (K)} & \textbf{Size} & \textbf{Mass ($M_\odot$)} & \textbf{Tracer} \\
		\midrule
		HI clouds & $\sim 100$ & tens of pc & --- & $21\,\mathrm{cm}$ line \\
		Diffuse/translucent molecular & $15\text{--}50$ & several pc & $3\text{--}100$ & CO, optical/UV absorption \\
		Giant Molecular Clouds (GMCs) & $\sim 10\text{--}20$ & $\sim 50\,\mathrm{pc}$ & $10^5\text{--}10^6$ & CO; sites of star formation \\
		Bok globules & $\sim 10$ & sub-pc & $1\text{--}1000$ & Dense, isolated dark cores \\
		\bottomrule
	\end{tabular}
	\caption{Approximate physical properties of interstellar cloud types (substantial overlap exists between categories).}
	\label{tab:ism_cloud_types}
\end{table}

The historical, purely observational term \emph{nebula} (Latin: ``cloud'') covers any extended, diffuse object, and is refined as follows: \emph{reflection nebulae} shine by starlight scattered off surrounding dust (e.g.\ the Pleiades); \emph{dark nebulae} are simply clouds of sufficiently high extinction to appear as silhouettes (GMCs, Bok globules); \emph{emission nebulae} radiate their own line emission, and are subdivided into $\mathrm{HII}$ regions (gas photoionized by nearby hot stars), planetary nebulae (ionized gas ejected during late stellar evolution, e.g.\ the Ring Nebula), and supernova remnants (shocked ejecta from a stellar explosion, e.g.\ the Crab Nebula, remnant of SN 1054, often hosting a central pulsar).

\section{Heating and Cooling of the ISM}
\label{sec:ism_heating_cooling}

The ISM's temperature is set by a local balance between heating and cooling processes, and this balance determines the thermal pressure that Section~\ref{sec:virial_theorem} will show is what resists gravitational collapse.

\paragraph{Heating.} Cosmic rays ($10\text{--}10^{14}\,\mathrm{MeV}$ particles accelerated in stellar flares and supernova shocks) ionize atoms and molecules; the ejected energetic electron then collisionally heats the surrounding gas. Ultraviolet photons from massive stars photoionize hydrogen and eject photoelectrons directly from dust-grain surfaces (the dominant heating mechanism in many regions), while stellar X-rays provide a harder, more penetrating ionizing source. Dust grains also absorb continuum radiation directly, and mechanical shocks from stellar winds and supernovae thermalize bulk kinetic energy.

\paragraph{Cooling.} Collisionally excited molecules and fine-structure atomic levels decay radiatively in the infrared; because infrared photons are scattered comparatively weakly by dust (Table~\ref{tab:scattering_regimes}, Rayleigh regime), they escape the cloud efficiently and carry thermal energy away with them. Thermal continuum emission from dust grains provides a second, generally slower, cooling channel.

Efficient cooling is a physical prerequisite for star formation: a cloud that could not radiate away the heat generated by compression would simply be pushed back by its own rising pressure. The next section makes this competition between gravity and pressure quantitative.

\section{Gravitational Instability: The Virial Theorem}
\label{sec:virial_theorem}

We now turn from radiative properties to self-gravitating dynamics: under what conditions does an interstellar cloud, supported until now by its own thermal (and bulk) pressure, instead collapse under its own weight? The essential tool is the \emph{virial theorem}, derived here from Newtonian mechanics applied to a self-gravitating fluid.

\subsection{The Second Moment of the Mass Distribution}

Consider a cloud of gas of density $\rho(\vecr, t)$ occupying a finite volume, isolated from external pressure. Define its (scalar) second moment of mass about the origin,
\begin{equation}
	I \equiv \int \rho\,r^2\,dV = \int r^2\,dm,
	\label{eq:moment_of_inertia_def}
\end{equation}
where the second form treats the integral as running over fixed fluid mass elements $dm = \rho\,dV$ following the flow (a Lagrangian sum). Differentiating once,
\begin{equation}
	\dv{I}{t} = \int 2\vecr\cdot\vec v\,dm,
	\label{eq:moment_first_derivative}
\end{equation}
and differentiating again,
\begin{equation}
	\dv[2]{I}{t} = 2\int v^2\,dm + 2\int \vecr\cdot\dv{\vec v}{t}\,dm = 4K_{\rm bulk} + 2\int \vecr\cdot\dv{\vec v}{t}\,dm,
	\label{eq:moment_second_derivative}
\end{equation}
where $K_{\rm bulk} \equiv \frac{1}{2}\int v^2\,dm$ is the cloud's bulk (macroscopic) kinetic energy.

\subsection{The Force Terms: Pressure and Self-Gravity}

Each fluid element obeys Euler's equation under pressure and self-gravity,
\begin{equation}
	\rho\,\dv{\vec v}{t} = -\nabla P - \rho\,\nabla\Phi,
	\label{eq:euler_equation}
\end{equation}
where $\Phi$ is the gravitational potential sourced by the cloud's own density via Poisson's equation, $\nabla^2\Phi = 4\pi G\rho$. Substituting into the last term of Eq.~\eqref{eq:moment_second_derivative}:
\begin{equation}
	\int \vecr\cdot\dv{\vec v}{t}\,dm = -\int \vecr\cdot\nabla P\,dV - \int \rho\,\vecr\cdot\nabla\Phi\,dV.
	\label{eq:force_term_split}
\end{equation}

\paragraph{Pressure term.} Integrating by parts and discarding the surface term (zero external pressure at the cloud's edge):
\begin{equation}
	-\int \vecr\cdot\nabla P\,dV = -\oint P\,(\vecr\cdot\hat{\vec n})\,dA + \int P\,(\nabla\cdot\vecr)\,dV = 3\int P\,dV,
	\label{eq:pressure_term_integration}
\end{equation}
using $\nabla\cdot\vecr = 3$ in three dimensions.

\paragraph{Self-gravity term.} A standard identity for a potential sourced entirely by its own density (proved by symmetrizing the double integral $\Phi(\vecr) = -G\int \rho(\vecr')/|\vecr - \vecr'|\,d^3r'$ under $\vecr \leftrightarrow \vecr'$) gives
\begin{equation}
	-\int \rho\,\vecr\cdot\nabla\Phi\,dV = U, \qquad U \equiv -\frac{1}{2}\iint \frac{G\rho(\vecr)\rho(\vecr')}{|\vecr - \vecr'|}\,d^3r\,d^3r',
	\label{eq:self_gravity_identity}
\end{equation}
where $U < 0$ is the cloud's total gravitational self-energy.

\subsection{The Scalar Virial Theorem}

Collecting Eqs.~\eqref{eq:moment_second_derivative}, \eqref{eq:pressure_term_integration}, and \eqref{eq:self_gravity_identity}:
\begin{equation}
	\dv[2]{I}{t} = 4K_{\rm bulk} + 6\int P\,dV + 2U.
	\label{eq:virial_pre_ideal_gas}
\end{equation}
For a monatomic ideal gas, the thermal energy density is $\frac{3}{2}P$, so the total thermal kinetic energy is $K_{\rm th} = \frac{3}{2}\int P\,dV$, i.e.\ $\int P\,dV = \frac{2}{3}K_{\rm th}$. Substituting and defining the total kinetic energy $K \equiv K_{\rm bulk} + K_{\rm th}$:
\begin{keyresult}
	\begin{equation}
		\boxed{\dv[2]{I}{t} = 4K + 2U} \qquad\Longleftrightarrow\qquad \frac{1}{2}\dv[2]{I}{t} = 2K + U.
		\label{eq:scalar_virial_theorem}
	\end{equation}
\end{keyresult}
This is the \textbf{scalar virial theorem} for a self-gravitating, thermally supported gas cloud.

\subsection{Virial Equilibrium and the Onset of Collapse}

A cloud in a long-term statistically steady state has, by definition, a bounded, non-growing $I(t)$, so its time-average satisfies $\avg{\ddot I} = 0$:
\begin{equation}
	2K + U = 0 \qquad (\text{virial equilibrium}).
	\label{eq:virial_equilibrium}
\end{equation}
More generally, Eq.~\eqref{eq:scalar_virial_theorem} is a statement about stability:
\begin{itemize}
	\item If $2K > |U|$: $\ddot I > 0$, the cloud's mean-square radius accelerates outward — the cloud disperses.
	\item If $2K < |U|$: $\ddot I < 0$ — self-gravity overwhelms thermal support and the cloud accelerates \emph{inward}: gravitational collapse.
\end{itemize}
The marginal case $2K = |U|$ defines the critical condition for instability, evaluated explicitly below.

\subsection{The Jeans Mass and Jeans Length}

Model the cloud as a uniform sphere of mass $M$, radius $R$, and temperature $T$. Its gravitational self-energy is computed by imagining the sphere assembled shell by shell, bringing successive layers of matter in from infinity. When a partially-built sphere has radius $r < R$, its already-accreted mass (uniform density) is $m(r) = M(r/R)^3$. Bringing in the next infinitesimal shell $dm = 4\pi r^2\rho\,dr = 3Mr^2/R^3\,dr$ from infinity releases gravitational energy $dU = -Gm(r)\,dm/r$; integrating over the full assembly from $r=0$ to $r=R$:
\begin{align}
	U &= -\int_0^R \frac{G\,m(r)}{r}\,dm = -\int_0^R \frac{GM(r/R)^3}{r}\cdot\frac{3Mr^2}{R^3}\,dr = -\frac{3GM^2}{R^6}\int_0^R r^4\,dr \notag \\
	&= -\frac{3GM^2}{R^6}\cdot\frac{R^5}{5} = -\frac{3}{5}\frac{GM^2}{R}.
	\label{eq:uniform_sphere_self_energy}
\end{align}
Its thermal kinetic energy, with $N = M/(\mu m_{\rm H})$ particles of mean molecular weight $\mu$, is $K = \frac{3}{2}Nk_BT = \frac{3}{2}\frac{M}{\mu m_{\rm H}}k_BT$. Imposing the marginal virial condition $2K = |U|$:
\begin{equation}
	3\,\frac{Mk_BT}{\mu m_{\rm H}} = \frac{3}{5}\frac{GM^2}{R} \quad\Longrightarrow\quad R = \frac{GM\mu m_{\rm H}}{5k_BT}.
	\label{eq:jeans_critical_radius}
\end{equation}
Eliminating $R$ in favor of the ambient density via $M = \frac{4}{3}\pi R^3\rho$ yields the \textbf{Jeans mass}, and eliminating $M$ instead yields the \textbf{Jeans length}:
\begin{keyresult}
	\begin{align}
		M_J &= \left(\frac{5k_BT}{G\mu m_{\rm H}}\right)^{3/2}\left(\frac{3}{4\pi\rho}\right)^{1/2}, \label{eq:jeans_mass} \\
		R_J &= \left(\frac{15\,k_BT}{4\pi G\mu m_{\rm H}\rho}\right)^{1/2}. \label{eq:jeans_length}
	\end{align}
\end{keyresult}
A cloud with $M > M_J$ (equivalently $R > R_J$ at fixed density) has insufficient thermal pressure support and is gravitationally unstable. Both quantities decrease with density and increase with temperature: hot, tenuous gas is stable; cold, dense gas is not — precisely why star formation is observed in cold ($T \sim 10\text{--}20\,\mathrm{K}$) molecular clouds (Table~\ref{tab:ism_cloud_types}) rather than in the warm diffuse ISM.

\subsection{A Dynamical Cross-Check: The Jeans Dispersion Relation}

The energy argument above is global and geometry-dependent (it assumed a uniform sphere); it is worth confirming the same physics from a local, dynamical perspective. Consider an infinite, isothermal (sound speed $c_s^2 = k_BT/\mu m_{\rm H}$), self-gravitating fluid at rest with uniform background density $\rho_0$, and perturb it slightly: $\rho = \rho_0 + \rho_1$, velocity $\vec v = \vec v_1$, potential $\Phi = \Phi_0 + \Phi_1$. To linear order in the perturbation, continuity, Euler (with $\delta P = c_s^2\rho_1$), and Poisson's equation read
\begin{align}
	\pdv{\rho_1}{t} + \rho_0\nabla\cdot\vec v_1 &= 0, \label{eq:jeans_lin_continuity} \\
	\rho_0\pdv{\vec v_1}{t} &= -c_s^2\nabla\rho_1 - \rho_0\nabla\Phi_1, \label{eq:jeans_lin_euler} \\
	\nabla^2\Phi_1 &= 4\pi G\rho_1. \label{eq:jeans_lin_poisson}
\end{align}
Seeking plane-wave solutions $\propto \exp[i(\vec k\cdot\vecr - \omega t)]$ (with $\vec v_1 = v_1\hat{\vec k}$ along the wavevector) turns the derivatives into algebraic factors, $\partial_t \to -i\omega$ and $\nabla \to i\vec k$:
\begin{equation}
	-i\omega\rho_1 + i\rho_0 k v_1 = 0, \qquad
	-i\omega\rho_0 v_1 = -ic_s^2k\rho_1 - i\rho_0k\Phi_1, \qquad
	-k^2\Phi_1 = 4\pi G\rho_1.
\end{equation}
The first equation gives $v_1 = \omega\rho_1/(\rho_0 k)$; the third gives $\Phi_1 = -4\pi G\rho_1/k^2$. Dividing the second equation by $-i$ gives $\omega\rho_0 v_1 = c_s^2k\rho_1 + \rho_0k\Phi_1$; substituting both $v_1$ and $\Phi_1$:
\begin{equation}
	\rho_0\omega\cdot\frac{\omega\rho_1}{\rho_0 k} = c_s^2 k\rho_1 + \rho_0 k\left(-\frac{4\pi G\rho_1}{k^2}\right) \quad\Longrightarrow\quad \frac{\omega^2}{k}\rho_1 = c_s^2 k\rho_1 - \frac{4\pi G\rho_0}{k}\rho_1.
\end{equation}
Dividing through by $\rho_1/k$ yields the \textbf{Jeans dispersion relation}:
\begin{keyresult}
	\begin{equation}
		\omega^2 = c_s^2k^2 - 4\pi G\rho_0.
		\label{eq:jeans_dispersion_relation}
	\end{equation}
\end{keyresult}
Short-wavelength (large-$k$) perturbations have $\omega^2 > 0$ and oscillate as ordinary sound waves, stabilized by pressure. Long-wavelength perturbations with $k$ below the critical wavenumber $k_J \equiv \sqrt{4\pi G\rho_0}/c_s$ have $\omega^2 < 0$: $\omega$ is purely imaginary, and the perturbation grows exponentially rather than oscillating — gravitational instability. The corresponding critical wavelength,
\begin{equation}
	\lambda_J = \frac{2\pi}{k_J} = c_s\sqrt{\frac{\pi}{G\rho_0}},
	\label{eq:jeans_dispersion_wavelength}
\end{equation}
scales with $\rho$ and $T$ in exactly the same way as the energy-method result, Eq.~\eqref{eq:jeans_length}, differing only in the $\mathcal O(1)$ numerical prefactor — as expected, since the two derivations model the cloud's geometry differently (an infinite uniform medium here, versus a finite uniform sphere above) but describe the same underlying competition between pressure and self-gravity.

\section{From Jeans Instability to Free-Fall Collapse}
\label{sec:freefall}

Once $M > M_J$, thermal pressure is, by construction, no longer sufficient to balance gravity. Discard pressure entirely and ask: how fast does the cloud collapse? In this pressure-free limit, each mass shell at radius $r$ falls solely under the gravity of the mass $M(r)$ interior to it — which, for a shell falling inward without crossing any other shell (homologous collapse), is exactly the classical two-body radial-infall problem of Chapter~6, with $M(r)$ playing the role of the central mass.

\paragraph{Free-fall as a degenerate Kepler orbit.} A particle released from rest at radius $R_0$ and falling under a fixed point mass $M$ traces the zero-angular-momentum limit of a bound Keplerian ellipse: as $L \to 0$ at fixed (negative) energy, the semi-latus rectum $\ell = a(1-e^2) \to 0$, so $e \to 1$ while $a$ remains finite — the ellipse degenerates into a straight line segment of length $2a = R_0$. This is \emph{not} the same $e=1$ limit as the marginally unbound parabolic trajectory of Chapter~6 (there, $e \to 1$ because $E \to 0$ at fixed $L$); here $e \to 1$ because $L \to 0$ at fixed $E < 0$. With this identification, $a = R_0/2$, and Kepler's equation, Eq.~\eqref{eq:kepler_eq}, together with the parametrization $r = a(1-\cos u)$, still applies. The particle starts at rest at apoapsis of the degenerate ellipse ($u=\pi$, $r = 2a = R_0$) and falls to the center ($u=0$, $r=0$). Using $\ell = u - \sin u$ and the mean motion $n = \sqrt{GM/a^3}$ (Kepler's Third Law, Eq.~\eqref{eq:kepler3}), the free-fall time is
\begin{equation}
	t_{\rm ff} = \frac{\ell(u=\pi) - \ell(u=0)}{n} = \frac{\pi}{n} = \pi\sqrt{\frac{a^3}{GM}} = \pi\sqrt{\frac{(R_0/2)^3}{GM}} = \frac{\pi}{2\sqrt2}\sqrt{\frac{R_0^3}{GM}}.
	\label{eq:freefall_time_kepler}
\end{equation}

\paragraph{Expressing $t_{\rm ff}$ in terms of density.} Writing $M = \frac{4}{3}\pi R_0^3\rho$ (the mass initially enclosed within $R_0$) and substituting into Eq.~\eqref{eq:freefall_time_kepler}, the radius $R_0$ cancels completely:
\begin{keyresult}
	\begin{equation}
		t_{\rm ff} = \sqrt{\frac{3\pi}{32\,G\rho}}.
		\label{eq:freefall_time_density}
	\end{equation}
\end{keyresult}
The free-fall time depends \emph{only} on the initial density, not on the size of the region considered — every shell of a uniform-density cloud, regardless of its radius, reaches the center at the same instant. This is precisely what justifies treating the collapse as homologous (shells do not cross) in the first place.

\begin{example}[Free-Fall Time of a Giant Molecular Cloud Core]
	Consider a dense GMC core of pure molecular hydrogen with number density $n_{\rm H_2} = 10^{4}\,\mathrm{cm^{-3}}$, so $\rho = n_{\rm H_2}\,(2m_{\rm H}) \approx 3.3 \times 10^{-17}\,\mathrm{kg/m^3}$. Then
	\[
	t_{\rm ff} = \sqrt{\frac{3\pi}{32\,(6.674\times10^{-11})(3.3\times10^{-17})}} \approx 1.4\times10^{13}\,\mathrm{s} \approx 4.5\times10^{5}\,\mathrm{yr}.
	\]
	This sub-million-year timescale — set purely by density, via Eq.~\eqref{eq:freefall_time_density} — is the characteristic dynamical timescale of star formation within a dense molecular clump, dramatically shorter than the $\sim 10^{7}$--$10^{8}\,\mathrm{yr}$ lifetime of the diffuse GMC that hosts it.
\end{example}

\section{Beyond Free-Fall: Fragmentation and the First Hydrostatic Core}

Two effects complicate the idealized homologous collapse of Section~\ref{sec:freefall}, and together they set the stage for the birth of an individual star.

\paragraph{Fragmentation.} As long as the collapsing gas remains optically thin, it radiates away the heat generated by compression efficiently (Section~\ref{sec:ism_heating_cooling}) and collapses very nearly isothermally. Because the Jeans mass falls with density at fixed temperature ($M_J \propto \rho^{-1/2}$, Eq.~\eqref{eq:jeans_mass}), a collapsing isothermal cloud carries an ever-\emph{shrinking} local Jeans mass, so sub-regions that were stable at the cloud's initial density can become independently unstable as the collapse proceeds. The single GMC-scale instability of Section~\ref{sec:virial_theorem} therefore tends to fragment hierarchically into progressively smaller, independently collapsing clumps — the physical origin of stellar multiplicity and of clustered star formation, and (qualitatively) of the stellar initial mass function.

\paragraph{The Opacity Limit and the First Hydrostatic Core.} Fragmentation cannot continue indefinitely: once a clump's central density is high enough that it becomes optically thick to its own cooling radiation, compressional heating can no longer be radiated away, the collapse ceases to be isothermal, and the equation of state stiffens toward an adiabat ($P \propto \rho^\gamma$). At this point the innermost region halts its dynamical (free-fall) collapse and settles into a new, small hydrostatic object in approximate virial balance — Eq.~\eqref{eq:virial_equilibrium} is restored, but now with $K$ supplied by adiabatic compressional heating rather than by the fixed external temperature of Section~\ref{sec:ism_heating_cooling}. This \emph{first hydrostatic core}, of order a Jupiter mass, continues to grow by accreting the still-infalling envelope, and marks the transition from ``interstellar cloud'' to ``protostar.''

The subsequent evolution of this object — its contraction toward the main sequence, and the nuclear ignition that finally halts the contraction permanently — is the domain of stellar structure and evolution, and lies beyond the scope of this observationally-oriented volume. It is, however, the direct dynamical continuation of the virial argument constructed in this chapter: the entire life of a star on the main sequence is itself nothing more than another instance of Eq.~\eqref{eq:virial_equilibrium}, now with thermal pressure sustained by nuclear burning rather than by residual gravitational compression.
